\section{Discussion}
For the hybrid stacked-voting ensemble studied here, the evidence supports three conclusions of different strength. First, the ensemble added no discrimination: its AUC was equivalent to the selected model's overall and on five datasets, and inconclusive, not inferior, on three. Second, \emph{under marginal feature removal} its SHAP explanations were less faithful: the effect was consistent in direction (seven of eight datasets, significant on three, never significantly reversed under NDF) and robust to the checks above, but varied in size, and the pooled decrease was significant only with the less conservative meta-analysis method. Third, the conclusion depends on the removal method: for untuned models the difference under conditional removal was small and not significant, whereas with tuned models a conditional deficit remained. Statements about faithfulness should therefore always name the removal method; ours do not establish that ensemble explanations are less faithful under every definition.

The component analysis refines the interpretation. The deficit is not a generic ``linear versus tree'' effect: XGBoost was explained as faithfully as logistic regression, whereas random forest and nearest neighbors were not, and HE was explained about as faithfully as random forest. We hypothesize that averaging step-like and piecewise-constant components spreads attribution over correlated features and makes the ensemble respond less predictably to off-manifold inputs \cite{hooker_mentch,chen}; the one mechanism we tested directly, non-additivity, was not supported.

\emph{Clinical relevance.} On diabetic retinopathy, HE's explanation was less faithful than the selected model's for \NshareDR\% of patients: a clinician reading HE's top-ranked feature would, on average, be pointed to a feature that accounts for less of the risk score. Studies proposing clinical ensembles should therefore compare against the selected single model, explain the fitted ensemble and state the removal method.

\emph{Limitations and future work.} We studied one hybrid architecture, binary tasks and public datasets of at most 5{,}000 patients, analyzed as case studies; Parkinsons has only 32 subjects; exclusions by the gap threshold were heavy for two datasets, although an exclusion-free measure gave the same pattern; logistic regression was the selected comparator in most untuned folds; and no clinicians evaluated the explanations. This paper covers the SHAP faithfulness part of a broader explanation-reliability framework. LIME explanations, their stability under small input changes, and agreement between explainers are left for future work, as are external validation and larger clinical cohorts.

\section{Conclusion}
On eight healthcare datasets, a hybrid stacked-voting ensemble matched the discrimination of the internally selected single model within a pre-specified margin overall, while its SHAP explanations were less faithful under marginal feature removal on seven datasets, significantly on three. The pooled decrease was significant with DerSimonian--Laird pooling but not with the more conservative Hartung--Knapp method, and it was smaller under conditional removal; it should be read as a consistent but not definitive deficit. Explanation faithfulness should be reported alongside accuracy, with the removal method stated, when ensembles are proposed for clinical use. The supplement (S1--S13), protocols, code and per-fold outputs are at \url{https://github.com/RiyaKamboj/xai-hybrid-ensemble-faithfulness}.
